\clearpage
\begingroup
\fancyhead[L]{\small THEORIE DER BEWEGUNG DES MONDES.}
\renewcommand{\headrulewidth}{0pt}
\SourceLeaf{630}
\noindent Nun muss hier aber der Coefficient von \(\cos M,\ =-e\) werden,
folglich \(\eta\eta=1-\tfrac32nn\) oder \(\eta=1-\tfrac34nn\).
Übrigens ist dieser Werth von \(\eta\) von dem wahren noch beträchtlich
verschieden, wovon die Ursache im Folgenden vorkommen wird. --- Übrigens
dürfen wir hier \(\mathit{Const.}\cos(v+\mathit{Const.})=0\) setzen,
da in dieser Integralgleichung ohnehin schon 2 arbiträre Grössen vorkommen
(\(e\) und \(M-\eta v\)), welche begreiflich so bestimmt werden können,
dass der Initialzustand genau dargestellt werde. Behielte man den Theil
\(\mathit{Const.}\cos(v+\mathit{Const.})\) bei, so würden bei der zweiten
Verbesserung Cirkelbögen in den Werth von \(p\) kommen, und bei der
3\textsuperscript{ten} sogar Quadrate von Kreisbögen u.~s.~f., obwohl bei
genauerer Untersuchung diese Kreisbögen eigentlich doch wieder nur periodische
Gleichungen geben und mit der Voraussetzung \(\mathit{Const.}=0\) wieder
zusammenfallen. Obige Bemerkung macht indessen diese Entwickelung unnöthig. ---
Wir haben also
\[
 h=AAk(1-\tfrac12nn-\tfrac34gg)
\]
\[
\begin{aligned}
 p={}&h\Bigl\{1-\tfrac14gg\cos2\lambda+nn\cos2E+\tfrac32nn\varepsilon\cos a\\
 &-\tfrac{15}{8}ne\cos(2E-M)+\tfrac58nne\cos(2E+M)-e\cos M\\
 &\qquad\Bigl[+\tfrac12nn\varepsilon\cos(2E+a)
                  -\tfrac72nn\varepsilon\cos(2E-a)\Bigr]\Bigr\}
\end{aligned}
\]
\[
\begin{aligned}
 \log p={}&\log h-\tfrac14ee-\tfrac14gg\cos2\lambda+nn\cos2E
                  +\tfrac32nn\varepsilon\cos a\\
 &-\tfrac{15}{8}ne\cos(2E-M)+\tfrac98nne\cos(2E+M)
                  -e\cos M-\tfrac14ee\cos2M\\
 &\qquad\Bigl[+\tfrac12nn\varepsilon\cos(2E+a)
                  -\tfrac72nn\varepsilon\cos(2E-a)\Bigr]
\end{aligned}
\]
\[
\begin{aligned}
 \log\frac{\mathrm{d}u}{\mathrm{d}v}
   ={}&-\tfrac32\log h-\tfrac12\log k+\tfrac14nn+\tfrac38gg+\tfrac12ee\\
 &-\tfrac{11}{4}nn\cos2E+\tfrac{15}{4}ne\cos(2E-M)
                   -\tfrac{13}{4}nne\cos(2E+M)\\
 &+\tfrac12gg\cos2\lambda+2e\cos M+\tfrac12ee\cos2M-3nn\varepsilon\cos a\\
 &\qquad\Bigl[-\tfrac{11}{8}nn\varepsilon\cos(2E+a)
                   +\tfrac{77}{8}nn\varepsilon\cos(2E-a)\Bigr]
\end{aligned}
\]
\[
\begin{aligned}
 &\frac{\mathrm{d}u}{\mathrm{d}v}h^{\frac32}k^{\frac12}=\\
 &(1+\tfrac14nn+\tfrac38gg+\tfrac32ee)
       \Bigl\{1-\tfrac{11}{4}nn\cos2E\\
 &\qquad+\tfrac{15}{4}ne\cos(2E-M)-6nne\cos(2E+M)\\
 &\qquad+\tfrac12gg\cos2\lambda+2e\cos M
                         +\tfrac32ee\cos2M-3nn\varepsilon\cos a\\
 &\qquad\Bigl[-\tfrac{11}{8}nn\varepsilon\cos(2E+a)
                         +\tfrac{77}{8}nn\varepsilon\cos(2E-a)\Bigr]\Bigr\}.
\end{aligned}
\]
\par\noindent\hspace*{\parindent}Folglich
\[
 h^3k=1+\tfrac12nn+\tfrac34gg+3ee
\]
\SourceLeaf{631}
\noindent und
\[
\begin{aligned}
 u={}&v-\tfrac{11}{8}nn\sin2E+\tfrac{15}{4}ne\sin(2E-M)-2nne\sin(2E+M)\\
 &+\tfrac14gg\sin2\lambda+2e\sin M+\tfrac34ee\sin2M-3n\varepsilon\sin a\\
 &\qquad\Bigl[-\tfrac{11}{16}nn\varepsilon\sin(2E+a)
                    +\tfrac{77}{16}nn\varepsilon\sin(2E-a)\Bigr].
\end{aligned}
\]
\par\noindent\hspace*{\parindent}Um die numerischen Werthe dieser ersten
Annäherung mit den wahren vergleichen zu können, schreiben wir folgende
Rechnungen her:
\[
\begin{array}{r@{\;}c@{\;}r@{\qquad}r@{\;}c@{\;}r}
 \log nn&=&7{,}747\,8192 &e&=&11340^{\prime\prime}\\
 \log1^{\prime\prime}&=&4{,}685\,5749 &\log e&=&4{,}054\,6131\\ \cline{2-3}
 &&3{,}062\,2443 &\log n&=&8{,}873\,9096\\ \cline{5-6}
 \log\frac{11}{8}&=&0{,}138\,3027 &&&2{,}928\,5227\\ \cline{2-3}
 &&3{,}200\,5470 &\log\frac{15}{4}&=&0{,}574\,0313\\ \cline{5-6}
 &&1586\overset{\prime\prime}{,}89 &&&3{,}502\,5540\\
 \multicolumn{2}{r}{\text{nach d. Beob.}}&1928^{\prime\prime}
                  &&&3180\overset{\prime\prime}{,}93\\
 &&&\multicolumn{2}{r}{\text{Beob.}}&4695^{\prime\prime}
\end{array}
\]
\[
\begin{array}{r@{\;}c@{\;}r@{\qquad}r@{\;}c@{\;}r}
 \log e&=&4{,}054\,6131 &\varepsilon&=&0{,}016\,802\\
 \log ee&=&8{,}109\,2262 &\log\varepsilon&=&8{,}225\,3610\\
 &&4{,}685\,5749 &\log n&=&8{,}873\,9096\\ \cline{2-3}\cline{5-6}
 &&2{,}794\,8011 &&&7{,}099\,2706\\
 \log\tfrac34&=&9{,}875\,0613 &\log3&=&0{,}477\,1213\\ \cline{2-3}\cline{5-6}
 &&2{,}669\,8624 &&&7{,}576\,3919\\
 &&467\overset{\prime\prime}{,}59 &&&4{,}685\,5749\\ \cline{5-6}
 &&&&&2{,}890\,8170\\
 &&&&&777\overset{\prime\prime}{,}71\\
 &&&\multicolumn{2}{r}{\text{Beob.}}&690^{\prime\prime}.
\end{array}
\]
\par\begin{center}\rule{.18\linewidth}{.4pt}\end{center}
\clearpage
\endgroup
